%% Related work text for "Motivation and significance" section
%% References: refs.bib (30 entries)

Text-to-SQL semantic parsing has advanced rapidly, driven by cross-domain benchmarks such as Spider~\cite{yu2018spider}, SParC~\cite{yu2019sparc}, and CoSQL~\cite{yu2019cosql}. Systems including RAT-SQL~\cite{wang2020ratsql}, NatSQL~\cite{gan2021natsql}, BRIDGE~\cite{lin2020bridge}, Seq2SQL~\cite{zhong2017seq2sql}, and DIN-SQL~\cite{pourreza2024dinsql} have steadily improved accuracy on relational databases, with recent work exploring pre-trained sequence-to-sequence models such as T5~\cite{raffel2020t5,xie2022t5sql} for text-to-SQL generation. However, these systems translate natural language into SQL queries designed for tabular relational schemas and cannot express IoT-specific operations such as anomaly detection, time-bucketed trend aggregation, or cross-collection comparison that are essential for sensor data analysis. Natural language interfaces to databases have a long history~\cite{androutsopoulos1995nlidb,li2014constructing}, and recent surveys~\cite{affolter2019comparative,katsogiannis2023survey} confirm that current NLIDB systems remain limited to SQL-based backends, leaving NoSQL time-series databases without comparable natural language access.

Large language models built on the Transformer architecture~\cite{vaswani2017attention} have demonstrated strong capabilities in code generation~\cite{chen2021evaluating,touvron2023llama,roziere2023codellama,guo2024deepseekcoder}, and models such as GPT-4~\cite{openai2023gpt4} and Qwen~\cite{bai2023qwen} can synthesize database queries from natural language. Yet LLM-only approaches suffer from hallucination and schema mismatch, and lack deterministic fallback when the model produces invalid queries. While constrained decoding methods such as PICARD~\cite{scholak2021picard} partially address structured output reliability, production deployments require a hybrid strategy that combines LLM flexibility with rule-based fallback.

IoT data management platforms~\cite{gubbi2013iot,alfuqaha2015iot} increasingly rely on time-series databases and NoSQL stores such as InfluxDB~\cite{influxdb} and MongoDB~\cite{mongodb} for high-velocity sensor streams in domains including indoor air quality monitoring~\cite{who2010iaq}. Monitoring tools like Grafana~\cite{grafana} provide dashboard-based visualization but require users to construct queries manually. No existing platform offers an end-to-end natural language pipeline from question to answer with automatic visualization for IoT sensor data. Text2Query addresses this gap through a four-stage pipeline with a QueryPlan intermediate representation that maps ten IoT-specific operation types---including statistical aggregation, anomaly detection, trend analysis, and cross-collection comparison---to MongoDB aggregation pipelines, enabling non-technical users to query sensor data without learning database query syntax.